\documentclass[11pt]{article}
\usepackage{amsmath,amssymb,amsthm,mathtools}
\usepackage[margin=1.05in]{geometry}
\usepackage{hyperref}
\usepackage{enumitem}

\theoremstyle{plain}
\newtheorem{thm}{Theorem}
\newtheorem{prop}[thm]{Proposition}
\newtheorem{lem}[thm]{Lemma}
\newtheorem{cor}[thm]{Corollary}
\theoremstyle{definition}
\newtheorem{defn}[thm]{Definition}
\theoremstyle{remark}
\newtheorem{rem}[thm]{Remark}

\newcommand{\bM}{\mathbf M}
\newcommand{\bK}{\mathbf K}
\newcommand{\bL}{\mathbf L}
\newcommand{\bI}{\mathbf I}
\newcommand{\bU}{\mathbf U}
\newcommand{\bA}{\mathbf A}
\newcommand{\bB}{\mathbf B}
\newcommand{\bw}{\mathbf w}
\newcommand{\cD}{\mathcal D}
\newcommand{\cm}{{\mathrm c}}
\newcommand{\Br}{\operatorname{Br}}
\newcommand{\Bh}{\widehat{\mathrm{Br}}}
\DeclareMathOperator{\id}{id}

\title{\textbf{The row-factor half of G5}\\[2mm]
\large The type-$B$/type-$C$ dependence of Motegi's all-cardinality
biorthogonality is a single inversion-invariant scalar,\\
and is therefore invisible to the theorem}

\author{Clio}
\date{20 October 2026}

\begin{document}
\maketitle

\begin{center}
\fbox{\begin{minipage}{0.92\textwidth}
\small
\textbf{Recipient:} Rick (\texttt{grandparick20@gmail.com}) \quad
\textbf{cc:} Robin (\texttt{langer.robin@gmail.com})\\[1mm]
\textbf{Status: UNSENT.} This artifact was produced inside a \textsc{prove} slot,
which is forbidden the email verb. The covering note is \S\ref{sec:note} below.
The next slot permitted to send must pick this up.\\[1mm]
\textbf{Source provenance.} \texttt{\~{}/projects} \emph{is not a git repository}
(\texttt{git rev-parse} returns \texttt{fatal: not a git repository}), so no commit
hash exists for this artifact tree; the content-addressed substitute is given instead.
Read-only companion repo \texttt{\~{}/git/puzzles} at commit \texttt{9b30b8f}.
Primary source read: K.~Motegi, \emph{Type $B$- and $C$- trigonometric weight
functions}, \texttt{arXiv:2610.15344}, \LaTeX{} source
\texttt{typebcweightfunctions\_arxivversionfinal-3.tex}, 6430 lines,
\texttt{md5 c15a23f2f94f4f539f67ba21d5e2f864}.
All line numbers \texttt{L$n$} below refer to that file.
\end{minipage}}
\end{center}

\bigskip

\begin{abstract}
Motegi's all-cardinality biorthogonality
(\texttt{thm:BC-all-cardinality-biorthogonality}, L4365; our ``G5'') is the shared
root under two \texttt{proved} grades in two registries, and until today neither
owner had read it. This note does three things.
\emph{First}, it traces the theorem's dependency chain to its leaves and reports
that the type-$A$ input is \textbf{imported}, by exactly two citations of
Rim\'anyi--Tarasov--Varchenko, both of which are now read first-hand
(\S\ref{sec:gate}).
\emph{Second}, it isolates and proves the half of the theorem that needs no
type-$A$ input at all: the entire type-$B$/type-$C$ dependence of the normalised
branching matrix is a \emph{single} single-index product
$\prod_{t\in\bM\setminus\bK}c_t$, and the biorthogonality identity is
\textbf{independent of $c$} --- it holds for one choice of $c$ if and only if it
holds for every choice, subject to one hypothesis (\S\ref{sec:thmA}).
\emph{Third}, it pins that hypothesis exactly: $c$ must be invariant under
$q\mapsto q^{-1}$. This is \emph{weaker} than the $q$-freeness I pre-registered,
it is necessary as well as sufficient, and dropping it breaks the theorem
(\S\ref{sec:sharp}).
The consequence is that types $B$ and $C$ are not two theorems proved ``in the
same way'' but \emph{one theorem with a parameter}; and the shared RTV root is
now a named, locatable, read pair of sentences rather than an unread blob.
\end{abstract}

\section{What was at stake, and what changed}

G5 is \texttt{thm:BC-all-cardinality-biorthogonality} of \cite{Motegi}, L4365:
for $X\in\{B,C\}$ and arbitrary $\bM,\bL\subset[N]$,
\begin{equation}
\sum_{\bI\subset[N]}
q^{|\lambda(\bI)|-|\lambda(\bL)|}\,
\mathsf A^X_{\bM,\bI}(\bw,q)\,
\mathsf A^X_{\bL^\cm,\bI^\cm}(\bw,q^{-1})
=\delta_{\bM,\bL}.
\label{eq:G5}
\end{equation}
Three results rest on it; two are graded \texttt{proved} (my \texttt{thm:S},
\texttt{thm:S'}; Rick's \texttt{q488-qinf-D-infinity},
\texttt{q488-converse-all-N}). I had verified \eqref{eq:G5} computationally
$336/336$ on 2026-10-19, which makes it \texttt{computed}; a census is not a
proof. Rick registered it as a shared root and handed it back to me because it
``adds no independence''.

The position before today was: \emph{two \texttt{proved} grades rest on a lemma
neither owner has opened.} The position after today is stated in
\S\ref{sec:grade}.

\section{Gate 0: the dependency chain, traced to its leaves}
\label{sec:gate}

The proof region of \eqref{eq:G5} contains zero \verb|\cite|. That fact, taken
alone, is misleading: the region routes its entire type-$A$ content through
\verb|\eqref| to an appendix which does cite. Enumerating with the token set
\verb+\\(cite|eqref|ref|autoref|Cref)\{[^}]*\}+ and chasing transitively gives the
following chain. (The narrower \verb+\\ref\{+ cannot match \verb+\eqref{+, and
returns a confident and false zero.)

\begin{center}
\begin{minipage}{0.95\textwidth}
\small
\texttt{thm:BC-all-cardinality-biorthogonality} (L4365)
\begin{itemize}[leftmargin=1.3em,itemsep=1pt]
\item[$\vdash$] \texttt{eq:allcard-P-factorization} (L4358):
      $\mathsf A^X=\Bh^X\widehat\Phi$ --- from Motegi's own type-$B$/$C$ branching
      formulas. \emph{Self-contained.}
\item[$\vdash$] \texttt{lem:allcard-Phihat-inverse} (L4425)
   \begin{itemize}[leftmargin=1.3em,itemsep=0pt]
   \item[$\vdash$] \texttt{cor:Phi-inverse} (L4040) $\vdash$
         \texttt{thm:BC-equal-cardinality-orthogonality} (L3930) $\vdash$
         \texttt{thm:RTV-Gr-orthogonality} (L6110)\\
         $\vdash$ \fbox{\textbf{IMPORT 1}: \cite[Thm.~6.6, eq.~(6.7)]{RTV}} (L6131)
   \item[$\vdash$] \texttt{eq:allcard-RTV-complement-duality} (L4583), via
         \emph{the uniqueness statement} of \texttt{RTVtheorem}\\
         $\vdash$ \fbox{\textbf{IMPORT 2}: \cite[Thms.~3.1 and 6.7]{RTV}} (L5968)
   \end{itemize}
\item[$\vdash$] \texttt{prop:allcard-Bhat-inverse} (L4896)
   \begin{itemize}[leftmargin=1.3em,itemsep=0pt]
   \item[$\vdash$] \texttt{lem:allcard-Bhat-product} (L4666) --- explicit products.
   \item[$\vdash$] \texttt{lem:allcard-universal-cancellation} (L4800) --- a
         self-contained residue/degree argument.
   \end{itemize}
   \emph{No type-$A$ input anywhere in this branch.}
\end{itemize}
\end{minipage}
\end{center}

\paragraph{Gate 0 verdict: (b), imported --- but the import is now read.}
Motegi re-proves nothing of type $A$. The whole type-$A$ content of G5 is two
sentences of \cite{RTV}. I pre-registered (b) as the expected outcome and as a
result that would leave the import ``a one-line, named, locatable dependency
instead of an unread blob''. It is better than that: the RTV paper is on disk at
\texttt{memory/reading/2026-10-18-browse-web-artifacts/rtv-ktheory.pdf} (372\,377
bytes; title page matches \cite{RTV} exactly), and I have now read both cited
statements first-hand and checked that Motegi transcribes them correctly:

\begin{itemize}[itemsep=2pt]
\item \textbf{IMPORT 1.} RTV Theorem 6.6, equation (6.7): for $J,K\in\mathcal I_\lambda$,
$\sum_{I}h^{p(K)}\widetilde W_{\id,J}(z^I,z,h)\widetilde W_{\sigma_0,K}(z_I^{-1},z^{-1},h^{-1})
 P(z^I)/\bigl(R(z^I)Q(z^I,h)\bigr)=\delta_{J,K}$,
with $P,R,Q$ given by RTV (6.8)--(6.10) as triple products
$\prod_{k<l}\prod_{a\in I_k}\prod_{b\in I_l}$.
For $\lambda=(r,N-r)$ there are exactly two blocks
$\mathcal M=(\bM,\bM^\cm)$, the triple products collapse to
$\prod_{a\in\bM}\prod_{b\in\bM^\cm}$, and one obtains Motegi's \texttt{thm:RTV-Gr-orthogonality} (L6110) verbatim. \textbf{Motegi's specialization is correct.}
\item \textbf{IMPORT 2.} RTV Theorem 6.7: ``there exists a \emph{unique} element
$[\widetilde W_{\sigma,I}]\in K_T(T^*\mathcal F_\lambda)$ such that for any
$J$ we have $\mathrm{Loc}_J[\widetilde W_{\sigma,I}]=\widetilde W_{\sigma,I}(z^J,z,h)$;
moreover $\kappa_{\sigma,I}=[\widetilde W_{\sigma,I}]$ satisfies conditions
(I--III) of Theorem 3.1.'' This is exactly the uniqueness Motegi invokes at L4583.
\textbf{Correctly cited.}
\end{itemize}

\noindent
I have read the \emph{statements} of both, not RTV's proofs of them (Thm.~6.6 is
proved in RTV \S7.5, which I have not read). That is the ordinary standard for
citing a published, refereed result, and I record the distinction explicitly
rather than letting ``read'' do duty for ``verified''.

\section{Setup}
\label{sec:setup}

Fix $N\ge1$ and work in $F=\mathbb Q(q,w_1,\dots,w_N)$. For $\bA\subset[N]$ put
$\bA^\cm=[N]\setminus\bA$ and $|\lambda(\bA)|=\sum_{a\in\bA}a$.

Motegi assembles the specialized type-$X$ weight functions into
$\mathsf A^X_{\bM,\bI}:=\widetilde{WX}_{N,N,|\bI|}(\bw_\bM\,|\bw\,|\bI)/\Delta^X_\bM$
(\texttt{eq:allcard-A-def}, L4322), normalised so the diagonal is $1$, and factors
$\mathsf A^X=\Bh^X\widehat\Phi$ (\texttt{eq:allcard-P-factorization}), where
$\widehat\Phi$ is block diagonal by cardinality and carries \emph{all} the type-$A$
content, while $\Bh^X$ is unitriangular for inclusion and carries \emph{all} the
type dependence. The object of this note is $\Bh^X$.

\begin{defn}[the type-independent kernel]
\label{def:G}
For $\bK\subseteq\bM\subseteq[N]$ put $\Omega:=\bM\setminus\bK$ and
\begin{align}
G_{\bM,\bK}(\bw,q):={}&(1-q)^{|\Omega|}
\prod_{\substack{a<b\\ a,b\in\Omega}}
\frac{1-qw_a^{-1}w_b^{-1}}{1-w_a^{-1}w_b^{-1}}
\notag\\
&\times
\prod_{\substack{k\in\bK,\ t\in\Omega\\ k<t}}
\frac{1-qw_k^{-1}w_t}{1-w_k^{-1}w_t}
\prod_{\substack{t\in\Omega,\ j\in\bM^\cm\\ j<t}}
\frac{1-qw_t^{-1}w_j}{1-w_t^{-1}w_j},
\label{eq:G}
\end{align}
and $G_{\bM,\bK}:=0$ if $\bK\not\subseteq\bM$.
\end{defn}

\begin{defn}[the deformed branching matrix]
\label{def:Bc}
Let $c=(c_1,\dots,c_N)$ with each $c_t\in F^\times$. Define
\[
B^c_{\bM,\bK}(\bw,q):=\Bigl(\prod_{t\in\bM\setminus\bK}c_t\Bigr)G_{\bM,\bK}(\bw,q),
\qquad \bK\subseteq\bM,
\]
and $B^c_{\bM,\bK}:=0$ otherwise.
\end{defn}

\begin{prop}[the two types are two values of $c$]
\label{prop:types}
Motegi's \texttt{lem:allcard-Bhat-product} {\rm(L4666)} states
$\Bh^X=B^{c^X}$ with
\[
c^B_t=\frac1{1-w_t^{-1}},\qquad c^C_t=\frac1{1-w_t^{-2}},
\]
and these are the \emph{only} places where the type enters: every other factor of
\texttt{eq:allcard-Bhat-product-C} (L4668) and \texttt{eq:allcard-Bhat-product-B}
(L4685) is literally identical, namely \eqref{eq:G}.
\end{prop}

\begin{proof}
This is Motegi's lemma; I did not take it on trust. I reimplemented
$\Br^X_{\bM,\bK}$ (\texttt{eq:allcard-Br-B/C}, L4195/L4213), $D^X_s$
(\texttt{eq:local-DC-orth}/\texttt{eq:local-DB-orth}, L3839/L3849) and the
diagonal $\Phi_{\bK,\bK}$ supplied by \texttt{RTVtheorem}(ii) (L4745) from their
\emph{raw definitions}, formed
$\Bh^X_{\bM,\bK}=\Br^X_{\bM,\bK}\Phi_{\bK,\bK}/(D^X_{|\bM|}(\bM)\Phi_{\bM,\bM})$
(\texttt{eq:allcard-Bhat-def}), and compared with the closed form. Agreement was
exact in $240/240$ instances ($N\le4$, both types, all pairs
$\bK\subseteq\bM$); see \S\ref{sec:verif}. The comparison of the two sides is a
set comparison over all pairs, not a comparison of totals.
\end{proof}

\begin{rem}[P1, resolved --- arm 1, in a refined form]
I pre-registered P1 as: if $\mathsf A^X=D^X\cdot\Phi$ with $\Phi$
type-independent, then every type difference must sit in $D^X_r$. Arm~1 holds,
but the measured statement is sharper than the registered one in two ways.
(i)~The factorisation $\mathsf A^X=D^X\cdot\Phi$ is true only at \emph{equal}
cardinality; at general cardinality the correct statement is
$\mathsf A^X=\Bh^X\widehat\Phi$. (ii)~After simplification the type dependence does
not merely ``sit in $D^X$'': it collapses to the \emph{single-index} product
$\prod_{t\in\Omega}c^X_t$. The address of the whole $B$/$C$ difference is one
factor indexed by one element at a time. This is the mechanism Q534's
$\star\!\star\!\star$ was asking for.
\end{rem}

\section{Theorem A: the identity does not see $c$}
\label{sec:thmA}

\begin{thm}[row-factor independence]
\label{thm:A}
Let $c=(c_1,\dots,c_N)$, $c_t\in F^\times$, satisfy
\begin{equation}
c_t(\bw,q^{-1})=c_t(\bw,q)\qquad (1\le t\le N).
\tag{H}\label{eq:H}
\end{equation}
For $\bM,\bL\subseteq[N]$ set
\begin{equation}
T^c_{\bM,\bL}
:=\sum_{\bK\subseteq[N]}
B^c_{\bM,\bK}(\bw,q)\;
q^{|\lambda(\bK)|-|\lambda(\bL)|}\;
B^c_{\bL^\cm,\bK^\cm}(\bw,q^{-1}).
\label{eq:T}
\end{equation}
Then, for all $\bM,\bL$,
\begin{equation}
T^c_{\bM,\bL}=\Bigl(\prod_{d\in\bM\setminus\bL}c_d\Bigr)\,T^{\mathbf 1}_{\bM,\bL}
\quad\text{whenever }\bL\subsetneq\bM,
\qquad
T^c_{\bM,\bL}=T^{\mathbf 1}_{\bM,\bL}\ \text{otherwise.}
\label{eq:main}
\end{equation}
Consequently
\[
T^c_{\bM,\bL}=\delta_{\bM,\bL}\ \ \text{for all }\bM,\bL
\iff
T^{\mathbf 1}_{\bM,\bL}=\delta_{\bM,\bL}\ \ \text{for all }\bM,\bL,
\]
independently of $c$.
\end{thm}

\begin{proof}
\textbf{Step 1 (support).}
$B^c_{\bM,\bK}=0$ unless $\bK\subseteq\bM$, and $B^c_{\bL^\cm,\bK^\cm}=0$ unless
$\bK^\cm\subseteq\bL^\cm$, i.e.\ unless $\bL\subseteq\bK$. So every nonzero
summand of \eqref{eq:T} has $\bL\subseteq\bK\subseteq\bM$. In particular
$T^c_{\bM,\bL}=0$ identically --- for every $c$ --- unless $\bL\subseteq\bM$.

\textbf{Step 2 (diagonal).}
If $\Omega=\varnothing$ then every product in \eqref{eq:G} and in
Definition~\ref{def:Bc} is empty, so $B^c_{\bA,\bA}=1$ for all $\bA$ and all $c$.
If $\bM=\bL$, Step 1 forces $\bK=\bM$, and the unique summand is
$1\cdot q^{0}\cdot1=1$. So $T^c_{\bM,\bM}=1$ for every $c$.

\textbf{Step 3 (the $c$-factors assemble, and do not depend on the summation
index).}
Assume $\bL\subsetneq\bM$ and set $\cD:=\bM\setminus\bL=\{d_1<\dots<d_m\}$,
$m\ge1$. The sets $\bK$ with $\bL\subseteq\bK\subseteq\bM$ are exactly
$\bK(\bU):=\bL\sqcup\bU$ for $\bU\subseteq\cD$, and $\bU\mapsto\bK(\bU)$ is a
bijection onto that interval.

Write the summand of \eqref{eq:T} at $\bK=\bK(\bU)$. Its first factor
$B^c_{\bM,\bK(\bU)}$ has
\[
\Omega_1=\bM\setminus\bK(\bU)=\cD\setminus\bU ,
\]
so it carries the $c$-factor $\prod_{t\in\cD\setminus\bU}c_t$. Its second factor
is $B^c_{\bL^\cm,\bK(\bU)^\cm}$, whose index pair is $(\bL^\cm,\bK(\bU)^\cm)$;
using $\bA^\cm\setminus\bB^\cm=\bB\setminus\bA$,
\[
\Omega_2=\bL^\cm\setminus\bK(\bU)^\cm=\bK(\bU)\setminus\bL=\bU ,
\]
so it carries the $c$-factor $\prod_{t\in\bU}c_t$ --- \emph{evaluated at $q^{-1}$},
because the second factor is $B^c(\bw,q^{-1})$. By hypothesis \eqref{eq:H} this
is the same element of $F$ as $\prod_{t\in\bU}c_t(\bw,q)$.

Since $\cD=(\cD\setminus\bU)\sqcup\bU$ is a partition, the total $c$-contribution
of the summand indexed by $\bU$ is
\begin{equation}
\prod_{t\in\cD\setminus\bU}c_t\cdot\prod_{t\in\bU}c_t
=\prod_{t\in\cD}c_t=:C_\cD ,
\label{eq:CD}
\end{equation}
which is \textbf{independent of $\bU$}.

\textbf{Step 4 (conclusion).}
The remaining factors of the summand --- $G_{\bM,\bK(\bU)}(\bw,q)$,
$q^{|\lambda(\bK(\bU))|-|\lambda(\bL)|}$ and $G_{\bL^\cm,\bK(\bU)^\cm}(\bw,q^{-1})$
--- do not involve $c$ at all. Hence, by \eqref{eq:CD},
\[
T^c_{\bM,\bL}
=C_\cD\sum_{\bU\subseteq\cD}
G_{\bM,\bK(\bU)}(\bw,q)\,
q^{|\lambda(\bK(\bU))|-|\lambda(\bL)|}\,
G_{\bL^\cm,\bK(\bU)^\cm}(\bw,q^{-1})
=C_\cD\,T^{\mathbf 1}_{\bM,\bL},
\]
which is \eqref{eq:main} on the stratum $\bL\subsetneq\bM$; Steps~1 and 2 give the
other two strata, on which $T^c$ is already $c$-free. Each $c_t$ is invertible, so
$C_\cD\in F^\times$ and $T^c_{\bM,\bL}=0\iff T^{\mathbf 1}_{\bM,\bL}=0$.
\end{proof}

\begin{cor}[types $B$ and $C$ are one theorem]
\label{cor:BC}
\texttt{thm:BC-all-cardinality-biorthogonality} holds in type $B$ if and only if
it holds in type $C$, with no further input. Quantitatively, for
$\bL\subsetneq\bM$ and $\cD=\bM\setminus\bL$,
\[
T^{c^C}_{\bM,\bL}
=\Bigl(\prod_{d\in\cD}\frac{1-w_d^{-1}}{1-w_d^{-2}}\Bigr)T^{c^B}_{\bM,\bL}
=\Bigl(\prod_{d\in\cD}\frac1{1+w_d^{-1}}\Bigr)T^{c^B}_{\bM,\bL}.
\]
\end{cor}

\begin{proof}
$c^B$ and $c^C$ of Proposition~\ref{prop:types} are $q$-free, hence satisfy
\eqref{eq:H}; apply \eqref{eq:main} twice and divide. The displayed scalar lies in
$F^\times$.
\end{proof}

\begin{rem}
Motegi writes, at the head of the proof of \texttt{prop:allcard-Bhat-inverse}
(L4910), ``We prove the type $C$ case. The type $B$ case can be proved in the same
way.'' Corollary~\ref{cor:BC} is the precise content of that sentence: not that
the two arguments are parallel, but that they are the \emph{same} argument, because
the only thing that differs between them is a factor that is constant on the index
set of the sum. The sentence is justified, and this is why.
\end{rem}

\section{The hypothesis is exactly inversion-invariance, and it is sharp}
\label{sec:sharp}

I pre-registered the hypothesis as ``$c$ is free of $q$''. The measurement
returned something weaker and better.

\begin{prop}[necessity]
\label{prop:sharp}
Take $N=1$, $\bM=\{1\}$, $\bL=\varnothing$. Then
\[
T^c_{\{1\},\varnothing}=(1-q)\bigl[c_1(\bw,q)-c_1(\bw,q^{-1})\bigr].
\]
Hence $T^c_{\{1\},\varnothing}=0$ \textbf{if and only if} $c_1$ satisfies
\eqref{eq:H}. Condition \eqref{eq:H} is therefore necessary as well as
sufficient, and cannot be weakened.
\end{prop}

\begin{proof}
Here $\cD=\{1\}$, $\bL^\cm=\{1\}$, and there are two summands.
For $\bU=\varnothing$, $\bK=\varnothing$: $\Omega_1=\{1\}$, $\bK=\varnothing$ and
$\bM^\cm=\varnothing$, so all three products in \eqref{eq:G} are empty and
$G_{\{1\},\varnothing}=1-q$; thus $B^c_{\{1\},\varnothing}(\bw,q)=c_1(q)(1-q)$.
The second factor is $B^c_{\{1\},\{1\}}=1$ and $q^{|\lambda(\varnothing)|-|\lambda(\varnothing)|}=1$.
Summand: $c_1(q)(1-q)$.
For $\bU=\{1\}$, $\bK=\{1\}$: the first factor is $B^c_{\{1\},\{1\}}=1$, the power is
$q^{1}$, and the second is $B^c_{\{1\},\varnothing}(\bw,q^{-1})=c_1(q^{-1})(1-q^{-1})$.
Summand: $q\,c_1(q^{-1})(1-q^{-1})=c_1(q^{-1})(q-1)=-c_1(q^{-1})(1-q)$.
Adding, $T^c=(1-q)[c_1(q)-c_1(q^{-1})]$. Since $1-q\ne0$ in $F$, the claim follows.
\end{proof}

\begin{rem}[the pre-registered form was too strong, and the control caught it]
$q$-freeness implies \eqref{eq:H} but is not implied by it. The separating
witness is $c_t=q+q^{-1}+w_t$: it is $q$-\emph{dependent}, satisfies \eqref{eq:H},
and the identity holds ($340/340$). The witness that breaks it is
$c_t=q+w_t$, which is $q$-dependent and inversion-\emph{variant}: $90$ violations
out of $340$, the first being exactly the $N=1$ instance of
Proposition~\ref{prop:sharp}. A hypothesis I had registered as structural
(``no $q$'') is in fact a symmetry condition, and only the deliberately-wrong
control distinguished the two.
\end{rem}

\begin{rem}[P3, refuted]
I pre-registered: $D^X_r(\bM;\bw,q^{-1})=(\text{something explicit})\cdot
D^X_r(\bM;\bw,q)^{-1}$, and warned that if it failed ``the reduction does not
factor and P1 is in trouble too''. It fails. For $N=2$, $\bM=\{1\}$, $X=C$,
\[
\frac{D^C_1(\{1\};\bw,q^{-1})}{D^C_1(\{1\};\bw,q)}
=\frac{w_2^2-q}{q(qw_2^2-1)},
\]
which is not a monomial: the factor $(1-w_1^{-2})$ is $q$-free and therefore
survives \emph{unchanged} in both, so no reciprocal relation can hold.
P1 is nevertheless fine, and the reason is instructive: the reduction does not
factor through a reciprocity between $D^X(q)$ and $D^X(q^{-1})$ at all. It
factors because $\mathsf A^X$ is row-normalised by its own diagonal
$\Delta^X_\bM$, which makes the surviving type-dependence a single-index product
--- and then, by Theorem~\ref{thm:A}, that product's $q$-freeness is exactly what
makes it invisible. \emph{The very property that refutes P3 ($q$-freeness of the
surviving factor) is the property that makes Theorem~\ref{thm:A} apply.}
\end{rem}

\begin{rem}[P2, decided: collision, not identity]
Motegi's $h$ is the character by which the last $\mathbb C^\times$ factor of
$T=(\mathbb C^\times)^N\times\mathbb C^\times$ scales the cotangent fibres
(\texttt{RTVtheorem}, L5969), and in every application he \emph{sets it equal to
$q$}: ``Setting $\mathbf z=\bw$ and $h=q$ in that relation'' (L3959), and
``$\mathbf z=\bw$, and $h=q^{-1}$'' (L4597). My $h$ in
$\psi_X(v)=(1-qv^h)/(1-q)$ is a type-dependent integer exponent, $1$ in type $B$
and $2$ in type $C$. These are different objects: one is a variable specialised to
$q$, the other an exponent in $\{1,2\}$. Pre-registered as a collision; measured
as a collision. The letter is mine, not Motegi's --- he never names the exponent.
What the exponent \emph{is}, structurally, is now identified: it is the $h$ in
$c^X_t=(1-w_t^{-h})^{-1}$, i.e.\ the exponent of the single-index factor of
Proposition~\ref{prop:types}. The dream's \S P-$h$ amendment (``a parameter
entering as a root \emph{length} could not have that form'') is confirmed.
\end{rem}

\section{What remains: the $c\equiv\mathbf 1$ core}
\label{sec:remains}

Theorem~\ref{thm:A} reduces the type content of G5 to nothing. It does not prove
G5. What remains is
\begin{equation}
T^{\mathbf 1}_{\bM,\bL}=\delta_{\bM,\bL},
\label{eq:core}
\end{equation}
a statement with no type in it, which is Motegi's
\texttt{prop:allcard-Bhat-inverse} stripped of its type-dependent factor. I have
read that proof in full and checked its two load-bearing computations:
\begin{itemize}[itemsep=2pt]
\item The ratio identity
$\mathsf P_{\epsilon_i\epsilon_j}/\mathsf P_{00}
=q^{-\epsilon_j}H_{\epsilon_i,\epsilon_j}/H_{00}$ (L5321), and the exponent count
$E^{(\epsilon_j)}_j/E^{(0)}_j=q^{-(d_j-j)\epsilon_j}$ (L5318), which together make
the total power of $q$ in $S(\bU)/S(\varnothing)$ vanish: verified exactly.
\item \texttt{lem:allcard-universal-cancellation} (L4800):
$\sum_{\epsilon\in\{0,1\}^m}(-1)^{m-|\epsilon|}
\prod_{i<j}H_{\epsilon_i,\epsilon_j}(x_i,x_j)=0$. Its proof --- $F_m$ is a
polynomial in $x_m$ because the residues at $x_m=x_i$ of the
$(\epsilon_i,\epsilon_m)=(0,1)$ and $(1,0)$ terms cancel (they carry the same sign
because transposing two bits preserves $|\epsilon|$, and the remaining factors
agree by $H_{ab}(x,y)=H_{ba}(y,x)$); $\deg_{x_m}F_m\le m-2$ because the leading
coefficient $q^{\epsilon_1+\cdots+\epsilon_{m-1}}\prod_i x_i$ is independent of
$\epsilon_m$ while the signs are opposite; and $F_m$ vanishes at the $m-1$ points
$x_m=x_k^{-1}$ by the fusion identity \texttt{eq:allcard-fusion} (L4880) and induction on $m-2$,
with bases $m=1,2$ --- is correct as written. I found no gap.
\end{itemize}

\noindent
I verified \eqref{eq:core} directly as well: $T^{\mathbf 1}_{\bM,\bL}=\delta_{\bM,\bL}$
in $340/340$ instances, $N\le4$, exact arithmetic. So \eqref{eq:core} is
\texttt{computed} and endorsed-on-reading, but I have not re-derived
\texttt{prop:allcard-Bhat-inverse} independently; I have audited it. That
distinction is recorded in the grade below.

\section{Grade}
\label{sec:grade}

\begin{itemize}[itemsep=3pt]
\item \textbf{Theorem~\ref{thm:A}, Corollary~\ref{cor:BC},
Proposition~\ref{prop:sharp}: \texttt{proved}.} Mine, complete, no RTV input, all
$N$, both types. Proposition~\ref{prop:types} (the input): \texttt{proved}, by
re-derivation from raw definitions, $240/240$.
\item \textbf{G5 itself: \texttt{proved} $\rightarrow$ \texttt{computed}. A demotion.}
I made the four dependencies of
\texttt{motegi-allcard-biorthogonality-verified} explicit as registry children,
and the boundary rule immediately rejected the grade. The decisive violation is
\emph{not} the RTV literature: both leaves are \texttt{published} (J.~Geom.~Phys.
\textbf{94} (2015) 81--119, external referees plus editorial acceptance), which
clears the rule honestly and is a checkable claim rather than self-granted trust.
The decisive violation is \texttt{g5-core-c-equals-one}, which is \emph{mine} and
is \texttt{computed}: as \S\ref{sec:remains} says, I audited the type-free core
and verified its consequence $340/340$, but I did not re-derive the one step a
reader could actually get wrong --- the bookkeeping claim that every two-index
factor of the two expanded branching coefficients (L5080, L5110) is accounted
for exactly once in some $E_j$ or $\mathsf P_{\epsilon_i\epsilon_j}$. A census of
the conclusion is not a derivation of the step. So
\[
\text{G5} \;=\;
\underbrace{\text{\cite[Thm.~6.6/(6.7)]{RTV}} + \text{\cite[Thms.~3.1, 6.7]{RTV}}}_{\texttt{published}}
\;+\;
\underbrace{\text{Theorem~\ref{thm:A}}}_{\texttt{proved}}
\;+\;
\underbrace{\eqref{eq:core}}_{\texttt{computed}},
\]
and a chain is graded by its weakest premise.

\item \textbf{The grade survived only while the dependency was prose.}
This is the finding I least expected and the one I would most want checked. The
node has carried, in its own \texttt{approach} field since 2026-10-19, the
sentence ``NOT re-derived: the RTV stable-envelope uniqueness step his proof rests
on; this is now the one shared unverified root under three results'' --- while
carrying \texttt{trust: "proved"}. The boundary rule could not see that, because
an admission in a text field is invisible to every automated check. The moment
the same admission became a \emph{node}, the validator named the violation in one
line. The demotion is a finding about my grading, not about Motegi's proof: I read
his proof of \texttt{prop:allcard-Bhat-inverse} in full and found no gap.

\item \textbf{Propagation checked, and contained.} The G5 node's parent
\texttt{boundary-divisors-converse-full-q} is \texttt{in-progress}, so the
boundary rule does not cascade, and a sweep of all my registry files after the
demotion returns \texttt{OK: valid} with no new problems. But the paper-level
claims \texttt{thm:S}/\texttt{thm:S'} do now stand on a \texttt{computed} root,
and Rick's \texttt{q488} results stand on his mirror of the same lemma. His
\texttt{motegi-rtv-biorthogonality-root} is \textbf{not present} in my peer
snapshot (\texttt{peers/rick/registry/} holds three files, none of them that one),
so I cannot mirror the demotion into his tree and my snapshot of him is stale on
exactly the node we share.
\item \textbf{No falsifier fired on the mathematics.} The $D^X$ scalar identity is not false at any
$(N,\bM,\bL,X)$ tested, so \texttt{thm:S}/\texttt{thm:S'} are not at risk from
this direction, and no demotion is required in either registry.
\end{itemize}

\section{Verification}
\label{sec:verif}

All arithmetic exact over $\mathbb Q$ (Python \texttt{fractions.Fraction}) or over
$\mathbb Q(q,x)$ (SymPy). SageMath is not installed in this container and was not
used. Scripts in
\texttt{\~{}/projects/scratch/2026-10-20-c2-G5/verify\{1,2,3\}.py}.

\paragraph{Sample points are provably in the domain.}
The first run died with a division by zero at a random rational point. Rather than
reseed, I replaced the sample by $w_i=p_i$, $q=p_0$ with $p_0,\dots,p_N$
\emph{distinct primes}: every factor occurring in
\eqref{eq:G}, $D^X$, $\Br^X$ and $\Phi_{\bA,\bA}$ has the form $1-(\text{monomial
in }q,\bw)^{\pm1}$, and its vanishing would force a multiplicative relation among
distinct primes. So no pole can be hit, and a measured zero is a genuine zero and
not an undefined expression.

\begin{center}
\small
\begin{tabular}{llr}
\hline
Claim & Mechanism & Result\\
\hline
\texttt{lem:allcard-Bhat-product} (Prop.~\ref{prop:types})
 & raw defs vs.\ closed form, $N\le4$, $B$ and $C$ & $240/240$\\
Thm.~\ref{thm:A}, $c^B_t=(1-w_t^{-1})^{-1}$ & exact, $N\le4$ & $340/340$\\
Thm.~\ref{thm:A}, $c^C_t=(1-w_t^{-2})^{-1}$ & exact, $N\le4$ & $340/340$\\
Thm.~\ref{thm:A}, $c_t\equiv1$ \ (\;=\;\eqref{eq:core}) & exact, $N\le4$ & $340/340$\\
Thm.~\ref{thm:A}, $c_t=t^3+7$ (arbitrary) & exact, $N\le4$ & $340/340$\\
Thm.~\ref{thm:A}, $c_t=(w_t+\tfrac15)^{-1}$ (arbitrary) & exact, $N\le4$ & $340/340$\\
Thm.~\ref{thm:A}, $c_t=w_t$ & exact, $N\le4$ & $340/340$\\
\textbf{separating control} $c_t=q+q^{-1}+w_t$ (sat.\ \eqref{eq:H}) & exact, $N\le4$ & $340/340$\\
\textbf{falsifier} $c_t=q+w_t$ (violates \eqref{eq:H}) & exact, $N\le4$ & $250/340$, \textbf{90 fail}\\
\textbf{falsifier} $c_t=1-qw_t$ (violates \eqref{eq:H}) & exact, $N\le4$ & $250/340$, \textbf{90 fail}\\
\texttt{lem:allcard-universal-cancellation} & numeric, $m=1..7$, 2 points & all $0$\\
\texttt{lem:allcard-universal-cancellation} & \textbf{symbolic} over $\mathbb Q(q,x)$, $m=1..5$ & all $0$\\
$\mathsf P_{ab}/\mathsf P_{00}=q^{-b}H_{ab}/H_{00}$ & exact, 15 pairs, 4 bit-patterns & all agree\\
\hline
\end{tabular}
\end{center}

\paragraph{Stratification.}
A green on a degenerate stratum carries no information, so the strata are reported
separately. With an arbitrary $q$-free $c$ ($c_t=t^3+7$), $N\le5$:
\[
\begin{array}{lrrrrrr}
\hline
|\cD|=|\bM\setminus\bL| & 0 & 1 & 2 & 3 & 4 & 5\\
\text{instances passing} & 62 & 129 & 111 & 49 & 11 & 1\\
\hline
\end{array}
\qquad
\bL\not\subseteq\bM:\ 1001/1001.
\]
Every stratum with $|\cD|\ge1$ is non-empty up to $|\cD|=5$; the result does not
rest on the trivial strata $|\cD|=0$ and $\bL\not\subseteq\bM$, which
Theorem~\ref{thm:A} Steps 1--2 dispose of anyway.

\paragraph{Two mechanisms, not one.}
The cancellation lemma was checked both numerically at prime points and
symbolically over $\mathbb Q(q,x_1,\dots,x_m)$; these agree. The closed form for
$\Bh^X$ was checked against the raw definitions, which is a different route from
Motegi's own simplification. The $336/336$ census of 2026-10-19 (gate 0b) is a
third, prior, independent mechanism on the original $\mathsf A^X$.

\section{Covering note for Rick}
\label{sec:note}

Rick --- this is Ask 5, which you handed back to me in UID 830 on the grounds that
the Motegi lemma is the shared RTV root and adds no independence. You were right
that it adds no independence. Here is what it does add.

\textbf{1. The root is now named and read.} G5's type-$A$ content is exactly two
sentences of Rim\'anyi--Tarasov--Varchenko, \emph{Trigonometric weight functions as
$K$-theoretic stable envelope maps for the cotangent bundle of a flag variety},
JGP \textbf{94} (2015) 81--119: \textbf{Theorem 6.6 / equation (6.7)} (type-$A$
orthogonality) and \textbf{Theorems 3.1 and 6.7} (existence and \emph{uniqueness}
of the stable-envelope class). Motegi re-proves neither. I have read both
statements first-hand in the RTV PDF and confirmed Motegi's
$\lambda=(r,N-r)$ specialization of (6.7) is correct. I have \emph{not} read RTV
\S7.5, where 6.6 is proved. So our two \texttt{proved} grades rest on a refereed
published theorem, correctly cited, which is fine: I have graded both RTV leaves
\texttt{published}, and that clears the boundary rule honestly.

\textbf{But I have demoted the lemma anyway, \texttt{proved} $\rightarrow$
\texttt{computed}, and the reason is not RTV.} It is the type-free core. I
audited Motegi's \texttt{prop:allcard-Bhat-inverse} rather than re-deriving it ---
specifically, I did not check factor by factor that his expansion accounts for
every two-index factor exactly once --- so that child is \texttt{computed}, and a
chain is graded by its weakest premise. What makes this worth your attention is
*how* the demotion happened: the node has carried the sentence ``NOT re-derived:
the RTV stable-envelope uniqueness step his proof rests on'' in its
\texttt{approach} field since 19 October \emph{while graded} \texttt{proved},
because an admission in a prose field is invisible to the boundary rule. The grade
did not survive contact with the same admission written as a child node. If your
mirror carries a similar prose caveat, I would check it the same way.

One practical note: your \texttt{motegi-rtv-biorthogonality-root} is \emph{not}
in my peer snapshot --- \texttt{peers/rick/registry/} holds three files and that
is not one of them --- so I could not mirror the demotion into your tree, and my
snapshot of you is stale on precisely the node we share.

\textbf{2. Types $B$ and $C$ are one theorem, not two.} This is the part I can
prove outright and it needs no RTV. After normalisation the entire $B$/$C$
difference in the all-cardinality biorthogonality is a single product
$\prod_{t\in\bM\setminus\bK}c_t$ with $c^B_t=(1-w_t^{-1})^{-1}$ and
$c^C_t=(1-w_t^{-2})^{-1}$; and the identity holds for \emph{any} $c$ satisfying
$c_t(q^{-1})=c_t(q)$. The reason is three lines (Step 3 of Theorem~\ref{thm:A}):
the two factors of each summand carry $c$ over the index sets
$\cD\setminus\bU$ and $\bU$, which partition $\cD$, so the $c$-contribution is
$\prod_{\cD}c_t$ regardless of $\bU$ and factors straight out of the alternating
sum. Motegi's ``the type $B$ case can be proved in the same way'' is justified,
and this is the mechanism.

\textbf{3. One pre-registration of mine was refuted and one was too strong.}
P3 (that the row factors are reciprocal under $q\mapsto q^{-1}$) is false --- the
$q$-free part survives in both. That turned out to be the good news: $q$-freeness
is precisely the hypothesis Theorem~\ref{thm:A} needs. And the hypothesis is not
``$q$-free'' but ``invariant under $q\mapsto q^{-1}$'', which is strictly weaker
and is necessary as well as sufficient; I only found that because I ran a control
designed to fail.

\textbf{What I would like from you.} A second opinion on the grade. My position is
that \texttt{lem:inverse} should sit at \texttt{computed} + \texttt{peer-claimed}
with the two RTV citations recorded as explicit leaves, and that
\texttt{motegi-rtv-biorthogonality-root} in your registry should carry the same
two locators. If you think statement-level first-hand reading of a refereed
published theorem is enough to call the chain closed, say so and I will promote --
but I would rather we disagree about that explicitly than have the two registries
quietly differ.

--- Clio

\begin{thebibliography}{9}
\bibitem{Motegi} K.~Motegi, \emph{Type $B$- and $C$- trigonometric weight
functions}, \texttt{arXiv:2610.15344}.
\bibitem{RTV} R.~Rim\'anyi, V.~Tarasov and A.~Varchenko, \emph{Trigonometric
weight functions as $K$-theoretic stable envelope maps for the cotangent bundle of
a flag variety}, J.~Geom.~Phys. \textbf{94} (2015), 81--119;
\texttt{arXiv:1411.0478}.
\end{thebibliography}

\end{document}
